% =============================================================================
%                              MODULE OVERVIEW
% =============================================================================
\section*{Module Overview}
\addcontentsline{toc}{section}{Module Overview}

\subsection*{About This Module}

This module focuses on \textbf{multi-agent systems}---systems consisting of multiple autonomous entities with distributed information, computational capabilities, and potentially conflicting objectives. These systems can be:

\begin{itemize}
  \item \textbf{Cooperative}: sensor networks, teams of mobile robots in a warehouse
  \item \textbf{Competitive}: electronic marketplaces, resource/task allocation
  \item \textbf{Mixed}: collaborative platforms like Wikipedia or Stack Overflow
\end{itemize}

Agents may be entirely artificial, entirely human, or a mix of both.

\textbf{Module code:} 6CCSACCM \hfill \textbf{Semester:} 2 (AY 2025/26)

\bigseparator

\subsection*{Module Aims}

The module aims to:
\begin{itemize}
  \item Discuss foundational and recent research on various aspects of \textbf{multi-agent systems}.
  \item Place particular emphasis on \textbf{coordination strategies} and \textbf{fostering cooperative behaviour}.
  \item Focus on \textbf{networks} and the \textbf{internet economy}.
  \item Introduce students to the field and help identify \textbf{potential research directions}.
  \item Develop skills in \textbf{critical analysis of research papers} and \textbf{formulating research questions}.
\end{itemize}

\bigseparator

\subsection*{Learning Outcomes}

By the end of this module, you will be able to:

\begin{center}
\renewcommand{\arraystretch}{1.4}
\begin{tabular}{@{} c p{12cm} @{}}
\toprule
\textbf{Code} & \textbf{Learning Outcome} \\
\midrule
CCM1 & Use \textbf{game-theoretic, algorithmic and computational tools} to model and analyse multi-agent strategic behaviour and its outcomes. \\
CCM2 & Use \textbf{game-theoretic tools} to foster cooperation and stability in multi-agent systems. \\
CCM3 & Use \textbf{algorithmic mechanism design techniques} to model and design incentives in multi-agent interaction. \\
CCM4 & Apply \textbf{computational social choice techniques} to analyse multi-agent decision-making and improve the quality and stability of collective decisions. \\
\bottomrule
\end{tabular}
\end{center}

\bigseparator

\subsection*{Key Topics (Expected)}

Based on the learning outcomes, this module is likely to cover:

\begin{center}
\renewcommand{\arraystretch}{1.3}
\begin{tabular}{@{} l p{10cm} @{}}
\toprule
\textbf{Area} & \textbf{Topics} \\
\midrule
Game Theory & Strategic games, Nash equilibrium, coordination games, Prisoner's Dilemma, repeated games \\
Cooperation & Mechanisms for fostering cooperation, stability concepts, coalition formation \\
Mechanism Design & Incentive design, auctions, voting mechanisms, truthfulness \\
Social Choice & Voting rules, Arrow's theorem, manipulation, collective decision-making \\
Networks & Network effects, internet economy, platform design \\
\bottomrule
\end{tabular}
\end{center}

\bigseparator

\subsection*{Assessment and Feedback}

\textbf{Formative assessment opportunities:}
\begin{itemize}
  \item \textbf{Discussion exercises}: peer discussion and feedback
  \item \textbf{Tutorial questions}: regular problem sheets with feedback during sessions
  \item \textbf{Discussion forums}: post questions on KEATS
  \item \textbf{Office hours}: direct feedback opportunities
\end{itemize}

\TODO{} Add assessment structure (exam format, coursework) when details are provided.

\bigseparator
